\section{征服以前：三条走廊怎样被准备出来}

\subsection{寿春之后，北方可以把目光转向何处}

诸葛诞之败并没有替司马昭消除所有反对者，却改变了资源的排列。寿春是淮河
与江北道路交会的军镇；围城期间，魏方必须维持近畿到淮南的粮运，孙吴的援助
则必须穿过水陆节点才可能抵城。城破后，司马氏得到的不只是一次军功，也是一
条不再长期吸收中枢注意力的走廊。把这一点同二六三年伐蜀相连，才能理解何以
西南战争成为可能；不能反过来断言寿春的每一位居民都在替日后伐蜀作选择。

二六二年钟会在关陇经营，是把这个新的余地译成道路、情报和粮秣的过程。
《三国志·钟会传》使我们看见受命与进军的次序，却未留下可以加总的征发簿。
因此，读者不应在地图上替魏军画出一条毫无摩擦的粗箭头。较可靠的画面是：
关中和陇右的官吏、守军和运输者被纳入一项准备，剑阁前的正面军又必须决定
何时继续等待、何时另寻路径。钟会、邓艾和诸葛绪不是三根并排移动的箭；他们
掌握的道路、消息与风险不同，因而可以把同一场战役推向彼此并不相同的结果。

\subsection{阴平之后，成都没有“自动投降”}

邓艾越阴平的叙事常以险绝取胜吸引读者。险绝确实是关键：一条可勉强通过的
山路能够让原先围绕剑阁部署的守军失去预警时间。但险绝不是神话许可。人马、
武器和随行粮食必须越过，抵达平原后的军队也仍要取得新的消息与地方配合。
《三国志》各传和《资治通鉴》对行军细节的铺陈并不相同，故不能把某一段夸张
的困难或某一句劝降复原成无疑的现场。

成都的选择在于：正面主力仍被剑阁牵住，新的威胁已在绵竹方向出现，能够临时
集结的兵力和可传递的命令都有限。刘禅出降可以确认为最高政权的决定，不能
由此说蜀地所有人以同样速度接受魏。钟会与姜维此后的谋划、魏军内部的冲突和
成都兵变恰好说明，降表之后的军人、降众和地方关系仍有自己的计算。征服者
获得一座城，首先获得的是必须安置和重新解释的网络。

\subsection{从西陵到建业：吴的防线为何不能只守一处}

二七二年西陵之战让西晋看见孙吴在上游仍有很强的回击能力。步阐降晋给荆州
一处可利用的裂缝，陆抗则把江面、城池和援军重新接合。胜利不是吴拥有一条
永不失效的长江线，而是某个熟悉峡口、能够协调兵船与守军的网络暂时成功。
陆抗死后，西晋并非突然得到渡江钥匙；它得到的是一个更值得长期准备的判断：
若只从荆州强渡，吴仍可能集中抵抗，必须让益州的船队、荆州的军队和淮南的
压力在不同时间同时发生。

王濬在益州造船因此不是伐吴前的一项技术装饰。木材、工匠、河道、训练和沿江
补给要经历多年，才会使西部资源能随水而下。羊祜在襄阳的经营同样不是一场
战役的序幕，而是让边境招纳、屯戍和对岸信息可以持续积累。杜预接替时面对的
是已经存在而仍需协调的接口。官方叙事能让我们知道这些人受任、死亡和出兵的
顺序；它不能替每位船工、江岸居民或吴方守将说明他们在何日得知消息、何日
改变选择。

二七九年的多路进军正是这种不同准备速度的汇合。吴不能只把守军投到一处渡口，
因为上游舟师、荆州军和淮南压力会互相改变其调度；西晋也不能只凭一支舰队获胜，
因为各路若不能在各自走廊保持供给，分兵反会造成新的空隙。孙皓出降把最高
权属交给晋廷，却没有缩短建业县廷、港口、旧军和江南家庭改换日常关系所需的
时间。二八〇年是并立政权的终点，不是道路和文书从此失去距离的日期。
